\section{Problema}
Nos piden un programa P que, dado un conjunto de entrenamiento $(X_t, y_t)$, aprenda con técnicas de computación evolutiva una transformación del los atributos: síntesis de atributos o reducción de dimensionalidad. Esa transformación se aplica después al conjunto de validación: $X_v'=v(X_v)$. 

Si entrenamos un modelo arbitrario M con los atributos originales y medimos el error absoluto medio (MAE) en validación, obtenemos \textit{e}. Si entrenamos el mismo m odelo M con los atributos transformados y obtenemos \textit{e'}, el método funciona cuando \textit{e'}<\textit{e}.
\subsection{Planteamiento}
\begin{itemize}
    \item El modelo M es arbitrario. Por eso evaluamos siempre con varios modelos --- regresión lineal, KNN y árbol de decisión --- y no con uno solo. 

    \item La transformación se aprende solo con el entrenamiento, siendo la validación lo único que usamos para medir. 

    \item Buscamos un método generalizable que funcione igual en datasets distintos sin ajustes a mano. 
\end{itemize}

\subsection{Datos}
Partimos de dos datasets de regresión: Diabetes (442 muestras, 10 atributos) y California Housing (20640 muestras, 8 atributos). Para comprobar que el método es generalizable, se usan de manera adicional 10 datasets de OpenML --- la lista y sus tamaños los tenemos en una sección posterior.

\subsection{Protocolo común}
\begin{itemize}
    \item Usamos una partición con 70\% de entrenamiento y 30\% de validación con semillas distintas para cada unas de las repeticiones.

    \item La métrica va a ser el MAE en el conjunto de la validación. 

    \item Para poder comprobar datasets con escalas muy distintas, usamos el MAE relativo a los datos originales, y el MAE/MAE(crudo), promediado con la media geométrica. Si es menor que 1, significa que hemos mejorado. 

    \item Contrastes estadísticos, incluido el test de Wilcoxon por pares, aplicado sobre el cociente de errores de cada pareja (dataset, partición).
\end{itemize}